\chapter{Construction de portefeuille II}\label{ch:rs:portfolio-construction-ii}

Déplacer d'un point de base l'alpha d'une action dans un portefeuille long only de cinquante titres en moyenne--variance le fait bouger de 0.03~\% de son capital en moyenne, et de 0.33~\% au
pire. Tirer à nouveau le bruit de la prévision (une prévision tout aussi bonne, avec d'autres erreurs) et la \omterm{def:rs:vectorised-backtests:turnover}{rotation du portefeuille} atteint 24~\%. Sur huit ans, son
information sur les actions ne se manifeste pas~: après coûts, son ratio de Sharpe de 0.57 se situe à moins d'une erreur type de celui d'une simple détention du marché. Sur des jours plutôt
que des mois, le tableau s'inverse~: la manière dont un portefeuille bouge compte davantage que la direction vers laquelle il pointe. Avec trois signaux réels qui se dégradent à des vitesses différentes et des coûts de
transaction qui croissent avec le carré de la taille de la transaction, réoptimiser chaque jour vers le portefeuille moyenne--variance donne un ratio de Sharpe net de $-0.51$~; se diriger en partie
chaque jour vers le \omterm{def:rs:portfolio-construction-ii:aim}{portefeuille cible} (aim portfolio) de Gârleanu et Pedersen donne 4.06. Ce chapitre construit \texttt{firm.allocation}~: optimisation robuste, Black--Litterman,
\omterm{def:rs:portfolio-construction-ii:erc}{parité de risque} sous deux formes et le \omterm{def:rs:portfolio-construction-ii:aim}{portefeuille cible}.

\section{Optimisation robuste}

\begin{definition}[Optimisation robuste du portefeuille, ensemble d'incertitude]\label{def:rs:portfolio-construction-ii:robust}
L'\emph{optimisation robuste du portefeuille}\index{optimisation robuste du portefeuille} maximise le pire cas de l'objectif sur un \emph{ensemble
d'incertitude}\index{ensemble d'incertitude} de données. Pour un ellipsoïde d'alphas $\{\alpha + \Omega^{1/2}u : \lVert u\rVert \le \kappa\}$, le rendement espéré dans le pire cas est
$\alpha^\top w - \kappa\sqrt{w^\top\Omega w}$, et le problème est un programme conique du second ordre (Goldfarb et Iyengar).
\end{definition}

La pénalité $\kappa\sqrt{w^\top\Omega w}$ croît avec la taille du portefeuille le long des directions dans lesquelles l'alpha est le moins certain, qui sont les directions
que chérissait le maximiseur d'erreur du chapitre~25. Le solveur du chapitre (\texttt{robust\_mv}) n'a pas besoin de solveur conique~; il procède par majoration~: puisque $\sqrt x \le x/(2s) + s/2$
avec égalité en $x = s^2$, chaque étape résout un programme quadratique dans lequel la racine carrée est remplacée par une quadratique qui la touche au portefeuille courant, et la
fonction objectif croît à chaque étape (le \cref{lst:rs:portfolio-construction-ii:robust}).

Le premier portefeuille du chapitre détient les cinquante plus grands titres de l'univers du chapitre~25, en long only, entièrement investi, au plus 10~\% par titre, reconstruit aux 95
fins de mois avec l'alpha du chapitre~25 et le modèle de risque du chapitre~24 ($\gamma = 20$). La version robuste utilise un ellipsoïde de trois erreurs types de l'alpha
entre les titres. Un point de base ajouté à l'alpha d'une action déplace le portefeuille moyenne--variance classique de 0.029~\% du capital en moyenne et de 0.33~\% au plus~; le
portefeuille robuste de 0.028~\% en moyenne et de 0.12~\% au plus. La robustesse ne change pas la sensibilité moyenne d'un portefeuille à sa prévision~; elle supprime les plus grands sauts,
qui surviennent lorsque deux titres sont presque interchangeables pour l'optimiseur.

\section{Vues combinées à l'équilibre}

\begin{definition}[Rendements d'équilibre implicites, modèle de Black--Litterman]\label{def:rs:portfolio-construction-ii:bl}
Les \emph{rendements d'équilibre implicites}\index{rendements d'équilibre implicites} d'un portefeuille de marché $w_m$ sont $\Pi = \gamma\Sigma w_m$~: les rendements espérés pour
lesquels $w_m$ est l'optimum moyenne--variance. Le \emph{modèle de Black--Litterman}\index{modèle de Black--Litterman} (Black et Litterman, 1992) traite $\Pi$ comme une loi a priori
de covariance $\tau\Sigma$ et des vues $Pr = q$ d'incertitude $\Omega$, traitées comme des observations, et optimise sur la moyenne a posteriori
$\mu = [(\tau\Sigma)^{-1} + P^\top\Omega^{-1}P]^{-1}[(\tau\Sigma)^{-1}\Pi + P^\top\Omega^{-1}q]$ et de covariance $\Sigma + [(\tau\Sigma)^{-1} + P^\top\Omega^{-1}P]^{-1}$.
\end{definition}

Black--Litterman déplace la question de «~quels sont les rendements espérés~?~» vers «~de combien la prévision doit-elle m'écarter du marché~?~». Sans vues, le
portefeuille optimal est le marché~; avec une vue tenue pour certaine, la moyenne a posteriori coïncide exactement avec lui~; le portefeuille intermédiaire est un mélange. La version du chapitre prend
les pondérations par capitalisation comme $w_m$, l'alpha de chaque action comme une vue sur son rendement en excès de $\Pi$, $\tau = 0.05$ (la valeur usuelle) et des incertitudes sur les vues égales à
$\tau$ fois la variance de chaque action. Sa sensibilité à un point de base est de 0.031~\% en moyenne et de 0.11~\% au plus, et son portefeuille reste proche de celui du marché~: sa
volatilité réalisée est de 19.6~\%, celle du marché de 19.5~\%, tandis que celle de la moyenne--variance est de 15.8~\%.

\section{Parité de risque et budgétisation du risque}

\begin{definition}[Contribution au risque, budgétisation du risque, portefeuille à contributions au risque égales, parité de risque]\label{def:rs:portfolio-construction-ii:erc}
La \emph{contribution au risque}\index{contribution au risque} de la position $i$ est $w_i(\Sigma w)_i/(w^\top\Sigma w)$, sa part d'Euler de la variance (livre~6, chapitre~20)~;
les parts somment à un. La \emph{budgétisation du risque}\index{budgétisation du risque} choisit des pondérations long only dont les contributions au risque égalent des budgets donnés~; le \emph{portefeuille à contributions au risque
égales}\index{portefeuille à contributions au risque égales} attribue à chaque position la même part (Maillard, Roncalli et Teïletche). \emph{Parité de
risque}\index{parité de risque} est la pratique qui consiste à allouer entre classes d'actifs selon des contributions au risque égales, en général avec effet de levier.
\end{definition}

\begin{proposition}[Budgets de risque par la méthode de Newton]\label{prop:rs:portfolio-construction-ii:budget}
La fonction $f(y) = \frac12 y^\top\Sigma y - \sum_i b_i\ln y_i$ est strictement convexe sur $y > 0$ et son minimiseur vérifie $y_i(\Sigma y)_i = b_i$~; les pondérations
$w = y/\sum_i y_i$ ont donc pour \omterm{def:rs:portfolio-construction-ii:erc}{contributions au risque} $b$.
\end{proposition}

\begin{proof}
La matrice hessienne $\Sigma + \mathrm{diag}(b_i/y_i^2)$ est définie positive, et $f \to \infty$ au bord et à l'infini, donc il existe un unique minimiseur, où
$\Sigma y = b/y$ composante par composante. Les \omterm{def:rs:portfolio-construction-ii:erc}{contributions au risque} sont invariantes à un changement d'échelle de $y$, et $\sum_i y_i(\Sigma y)_i = \sum_i b_i$.
\end{proof}

Les portefeuilles fondés sur le risque ignorent l'alpha, si bien qu'ils ne bougent pas quand il change~: leur sensibilité est nulle, par construction et sans mérite. Leur vertu se trouve ailleurs.
Maillard, Roncalli et Teïletche ont montré que la volatilité du \omterm{def:rs:portfolio-construction-ii:erc}{portefeuille à contributions au risque égales} se situe entre celle du portefeuille de variance minimale et
celle du portefeuille à pondérations égales, et c'est le cas ici~: 18.4~\% contre 19.3~\% pour les pondérations égales. Asness, Frazzini et Pedersen ont expliqué pourquoi la \omterm{def:rs:portfolio-construction-ii:erc}{parité de risque} entre classes d'actifs a
rapporté davantage que son risque ne le laisse croire~: les investisseurs qui répugnent au levier font monter les actifs risqués, ce qui laisse aux actifs plus sûrs des rendements ajustés du risque plus élevés pour quiconque
peut recourir au levier sur ces actifs.

\section{Allocation hiérarchique}

\begin{definition}[Parité de risque hiérarchique]\label{def:rs:portfolio-construction-ii:hrp}
La \emph{parité de risque hiérarchique}\index{parité de risque hiérarchique} (HRP, López de Prado) ordonne les actifs par une classification hiérarchique à liaison simple des distances de corrélation
$\sqrt{(1-\rho_{ij})/2}$, et répartit le capital de façon récursive~: chaque bloc de la liste ordonnée est coupé en deux, et chaque moitié reçoit du capital en proportion inverse
à la variance de leurs portefeuilles de variance inverse.
\end{definition}

La \omterm{def:rs:portfolio-construction-ii:hrp}{parité de risque hiérarchique} n'inverse jamais la matrice de covariance, si bien qu'elle ne peut pas amplifier ses plus petites valeurs propres~; sur une covariance diagonale, elle donne les
pondérations de variance inverse, et sur une matrice en blocs elle répartit d'abord le capital entre les blocs avant de regarder à l'intérieur de chacun. Sur les cinquante titres, elle détient au plus
5.8~\% du capital dans un seul titre et réalise une volatilité de 17.8~\%. La comparaison des huit règles sur 95 mois, après des coûts de dix points de base par transaction~:

\begin{center}
\small
\setlength{\tabcolsep}{4pt}%
\begin{tabular}{@{}p{5.0cm}rrrrr@{}}
\toprule
règle & rotation au nouveau tirage & Sharpe & après coûts & volatilité & rotation\\
\midrule
moyenne--variance, modèle factoriel & 23.9~\% & 0.62 & 0.57 & 15.8~\% & 33~\%\\
moyenne--variance, covariance échantillonnale & 23.4~\% & 0.82 & 0.77 & 15.9~\% & 31~\%\\
moyenne--variance robuste & 22.4~\% & 0.66 & 0.61 & 15.8~\% & 31~\%\\
Black--Litterman & 24.4~\% & 0.75 & 0.72 & 19.6~\% & 27~\%\\
\omterm{def:rs:portfolio-construction-ii:erc}{contributions au risque} égales & 0 & 0.66 & 0.65 & 18.4~\% & 7~\%\\
\omterm{def:rs:portfolio-construction-ii:hrp}{parité de risque hiérarchique} & 0 & 0.69 & 0.67 & 17.8~\% & 14~\%\\
pondérations égales & 0 & 0.67 & 0.66 & 19.3~\% & 5~\%\\
pondérations par capitalisation & 0 & 0.66 & 0.65 & 19.5~\% & 6~\%\\
\bottomrule
\end{tabular}
\end{center}

La rotation au nouveau tirage est la rotation à sens unique lorsque le bruit de la prévision est tiré à nouveau~; la dernière colonne est la rotation mensuelle à sens unique. Un ratio de Sharpe mesuré sur 7.9 ans a une erreur type d'environ 0.4, et toutes les lignes sont à moins d'une erreur type les unes des autres. Les portefeuilles long only sont dominés par le
marché qu'ils détiennent~; leurs alphas ne sont qu'une légère inclinaison, et le bruit de la prévision, tiré à nouveau, déplace les portefeuilles fondés sur l'alpha d'un quart de leur capital. C'est là la
conclusion honnête d'une comparaison souvent présentée comme un concours~: les règles diffèrent par leur stabilité, leur concentration et leur coût, mais pas de façon détectable dans leur performance.

\begin{omfigure}
\begin{tikzpicture}
\begin{axis}[name=a,width=7cm,height=5.4cm,ybar,bar width=7pt,ylabel={\small rotation par nouveau tirage},xtick=data,symbolic x coords={MV,sample MV,robust,BL,ERC,
  HRP,equal,cap},xticklabels={MV,MV empirique,MV robuste,BL,ERC,HRP,équipondéré,capitalisation},x tick label style={rotate=45,anchor=east,font=\footnotesize},tick label style={font=\footnotesize},ymin=0,ymax=0.3,grid=major,
  grid style={gray!25},yticklabel style={/pgf/number format/fixed,/pgf/number format/precision=2},scaled y ticks=false,table/col sep=comma]
\addplot[fill=omDef!55,draw=omDef] table[x=label,y=redraw]{figdata/research/26-portfolio-construction-ii/methods.csv};
\end{axis}
\begin{axis}[at={(a.east)},anchor=west,xshift=1.2cm,width=7cm,height=5.4cm,ybar,bar width=7pt,ylabel={\small ratio de Sharpe après coûts},xtick=data,
  symbolic x coords={MV,sample MV,robust,BL,ERC,HRP,equal,cap},xticklabels={MV,MV empirique,MV robuste,BL,ERC,HRP,équipondéré,capitalisation},x tick label style={rotate=45,anchor=east,font=\footnotesize},tick label style={font=\footnotesize},
  ymin=0,ymax=1,grid=major,grid style={gray!25},table/col sep=comma]
\addplot[fill=gray!45,draw=gray] table[x=label,y=sr_net]{figdata/research/26-portfolio-construction-ii/methods.csv};
\end{axis}
\end{tikzpicture}

\omcaption{Les huit règles à positions longues uniquement sur cinquante titres~: à gauche, la rotation à sens unique lorsque le bruit de la prévision est tiré à nouveau~; à droite, le ratio de Sharpe après coûts
sur 95 mois (erreur type d'environ 0.4). Données~: \texttt{rs\_allocation}.}\label{fig:rs:portfolio-construction-ii:methods}
\end{omfigure}

\section{Optimisation multipériode avec coûts}

\begin{definition}[Portefeuille cible, optimisation multipériode]\label{def:rs:portfolio-construction-ii:aim}
L'\emph{optimisation multipériode}\index{optimisation multipériode} choisit les transactions en planifiant une suite de portefeuilles futurs sous des prévisions de rendements, de risques
et de coûts, et n'exécute que le premier (Boyd et ses coauteurs). Avec des coûts quadratiques $\frac\lambda2\Delta x^\top\Sigma\Delta x$ et des signaux dont le pouvoir
prédictif décroît aux taux $\phi_k$, Gârleanu et Pedersen montrent que le portefeuille optimal transige à chaque période une fraction $a/\lambda$ du chemin qui sépare le portefeuille courant du
\emph{portefeuille cible}\index{portefeuille cible} $(\gamma\Sigma)^{-1}\sum_k B_k f_k/(1 + \phi_k a/\gamma)$, qui pondère davantage les signaux lents que les signaux rapides, avec $a$ la
racine positive de $a^2(1-\rho) + a(\gamma(1-\rho) + \lambda\rho) - \gamma\lambda(1-\rho) = 0$ et $\rho$ le taux d'actualisation.
\end{definition}

Le second portefeuille du chapitre détient les mêmes cinquante titres en position longue et courte, chaque jour, avec les trois composantes de rendement espéré simulées comme signaux~: la
dérive persistante (demi-vie de 504 jours, $\phi = 0.00137$ par jour), la dérive post-annonce (soixante jours, $\phi = 0.0167$) et le retour à un jour ($\phi = 1$). Ils
sont les véritables rendements espérés, si bien que les ratios de Sharpe bruts sont élevés~; la question est seulement de savoir ce que coûte leur négociation. $\gamma = 72$ place le portefeuille markowitzien quotidien
à une volatilité annuelle de 10~\%. Quatre règles, à un niveau de coût $\lambda = 100$~:

\begin{center}
\small
\setlength{\tabcolsep}{4pt}%
\begin{tabular}{@{}p{5.6cm}rrrr@{}}
\toprule
règle, $\lambda = 100$ & ratio de Sharpe & après coûts & coûts annuels & rotation quotidienne\\
\midrule
réoptimisation quotidienne vers Markowitz & 5.14 & $-0.51$ & 56.6~\% & 2.11\\
\omterm{def:rs:portfolio-construction-ii:aim}{portefeuille cible} (transaction de 0.56 de l'écart) & 4.74 & 4.06 & 4.0~\% & 0.56\\
ajustement partiel vers Markowitz (0.56) & 5.07 & 3.36 & 12.6~\% & \\
rééquilibrage mensuel vers les signaux lents & 3.01 & 2.96 & 0.2~\% & \\
\bottomrule
\end{tabular}
\end{center}

Le portefeuille markowitzien quotidien poursuit le retour à un jour, le plus fort et le plus bref des trois signaux, et en paie le prix chaque jour. L'ajustement partiel transige moins mais
vise toujours le retour à un jour. Le \omterm{def:rs:portfolio-construction-ii:aim}{portefeuille cible} vise surtout les signaux lents~: avec $a/\gamma = 0.78$, le retour à un jour entre avec un poids de $1/(1 + 0.78) = 0.56$, la
dérive post-annonce avec 0.987 et la dérive persistante avec 0.999, et le portefeuille se déplace de 56~\% du chemin chaque jour. À dix fois le coût ($\lambda = 1\,000$)
la règle cible transige 0.23 de l'écart et dégage encore 3.15~; la réoptimisation quotidienne dégage $-22.62$ avec des coûts de 566~\% par an, l'ajustement partiel 1.08, le
rééquilibrage mensuel 2.55. À $\lambda = 10\,000$, seule la règle cible survit, à 1.61 (\cref{fig:rs:portfolio-construction-ii:lambda}).

\begin{omfigure}
\begin{tikzpicture}
\begin{semilogxaxis}[width=11cm,height=6cm,xlabel={niveau de coût $\lambda$},ylabel={ratio de Sharpe après coûts},tick label style={font=\small},label style={font=\small},
  xmin=8,xmax=12000,ymin=-3.6,ymax=5.2,grid=major,grid style={gray!25},legend columns=2,
  legend style={font=\footnotesize,at={(0.5,-0.25)},anchor=north,draw=none,fill=none,/tikz/every even column/.append style={column sep=8pt}},
  table/col sep=comma]
\addplot[black!60,thick,mark=o] table[x=lam,y=daily_Markowitz]{figdata/research/26-portfolio-construction-ii/lambda.csv};
\addplot[omDef,very thick,mark=*] table[x=lam,y=aim_portfolio]{figdata/research/26-portfolio-construction-ii/lambda.csv};
\addplot[omDef!55,thick,mark=square*] table[x=lam,y=partial_adjustment]{figdata/research/26-portfolio-construction-ii/lambda.csv};
\addplot[omThm,thick,dashed,mark=triangle*] table[x=lam,y=monthly_Markowitz]{figdata/research/26-portfolio-construction-ii/lambda.csv};
\draw[black!50,densely dotted] (axis cs:8,0) -- (axis cs:12000,0);
\legend{Markowitz quotidien,portefeuille cible,ajustement partiel,Markowitz mensuel}
\end{semilogxaxis}
\end{tikzpicture}

\omcaption{Ratio de Sharpe après coûts quadratiques en fonction du niveau de coût, pour quatre règles de transaction sur les trois signaux simulés~; les valeurs inférieures à $-3.5$ sont tracées à
$-3.5$. Données~: \texttt{rs\_allocation.multiperiod}.}\label{fig:rs:portfolio-construction-ii:lambda}
\end{omfigure}

\section{Tutoriel~: le basculement d'un point de base}\label{tut:rs:portfolio-construction-ii:tut}

\begin{tutorial}
\textbf{Objectif.} Construire le portefeuille de cinquante titres par huit règles, mesurer leur stabilité et leur performance réalisée, puis négocier quotidiennement les trois signaux simulés selon
quatre règles et une gamme de coûts. \textbf{État final~:} les deux tableaux, \cref{fig:rs:portfolio-construction-ii:methods,fig:rs:portfolio-construction-ii:lambda}.
\begin{tutsteps}
\item \textbf{Moyenne--variance robuste} par majorisation.
\omcode{code/firm/allocation/firm_allocation.py}{48}{59}{Une suite de programmes quadratiques pour l'objectif robuste.}\label{lst:rs:portfolio-construction-ii:robust}
\item \textbf{Budgets de risque} par la méthode de Newton.
\omcode{code/firm/allocation/firm_allocation.py}{85}{101}{Pondérations à contributions au risque données.}\label{lst:rs:portfolio-construction-ii:budget}
\item \textbf{Le \omterm{def:rs:portfolio-construction-ii:aim}{portefeuille cible} et le taux de transaction.}
\omcode{code/firm/allocation/firm_allocation.py}{157}{170}{La règle de Gârleanu et Pedersen.}\label{lst:rs:portfolio-construction-ii:gp}
\item \textbf{Exécuter} \texttt{rs\_allocation.stability}, \texttt{redraw\_turnover}, \texttt{backtest} pour chaque règle, \texttt{multiperiod} pour chaque règle et chaque coût, et
\texttt{fig\_allocation.py}.
\end{tutsteps}
\textbf{Pour aller plus loin.} Remplacer les vrais signaux par des prévisions bruitées et observer la progression de l'avantage de la règle cible (elle moyenne le bruit, comme la limite de rotation du chapitre~25
)~; donner au \omterm{def:rs:portfolio-construction-ii:erc}{portefeuille à contributions au risque égales} un effet de levier pour égaler la volatilité du portefeuille moyenne--variance, puis comparer.
\end{tutorial}

\section{Construction~: la boîte à outils d'allocation}\label{bld:rs:portfolio-construction-ii:allocation}

\begin{build}
\bfield{Rôle} Les alternatives à la simple moyenne--variance qu'utilise la société, chacune avec un test de sa propriété définissante, et une règle de transaction pour des signaux de
vitesses différentes.
\bfield{Interface} \texttt{mean\_variance(alpha, Sigma, gamma, lo, hi, budget)}, \texttt{robust\_mv(alpha, Sigma, Omega, kappa, gamma, lo, hi, budget)},
\texttt{implied\_returns(Sigma, w\_mkt, gamma)}, \texttt{black\_litterman(pi, Sigma, P, q, Omega, tau)}, \texttt{risk\_contributions(w, Sigma)},
\texttt{risk\_budget(Sigma, b)}, \texttt{erc(Sigma)}, \texttt{hrp(Sigma)}, \texttt{gp\_trade\_rate(gamma, lam, rho)}, \texttt{gp\_aim(Sigma, signals, phis, gamma, a)}.
\bfield{Règles} Chaque règle est jugée sur sa stabilité face à une prévision tirée à nouveau et sur sa performance après coûts avec son erreur type~; le taux de décroissance d'un signal est
estimé (chapitre~13) avant qu'il ne soit négocié~; les signaux rapides ne sont négociés que par une règle qui sait qu'ils sont rapides.
\bfield{Tests d'acceptation} \texttt{code/firm/allocation/tests/}~: la moyenne--variance face à sa forme fermée~; l'objectif robuste au-dessus de celui du portefeuille simple, et égal
à lui pour $\kappa = 0$~; les rendements implicites reproduisant le marché~; une vue certaine retenue et une vue sans valeur ignorée~; les \omterm{def:rs:portfolio-construction-ii:erc}{contributions au risque} égales à leurs budgets~; la HRP
égale à la variance inverse sur une matrice diagonale et répartissant également les blocs~; les limites du taux de transaction et la moindre pondération du \omterm{def:rs:portfolio-construction-ii:aim}{portefeuille cible}.
\bfield{Pour aller plus loin} La règle cible avec des coûts linéaires (régions de non-transaction)~; l'optimisation robuste sur la covariance également~; les budgets de risque hiérarchiques.
\end{build}

\begin{omsources}
\item F.~Black et R.~Litterman, ``Global portfolio optimization'', \emph{Financial Analysts Journal} 48(5), 1992~; PyPortfolioOpt, \texttt{black\_litterman}
(implémentation documentée, d'après He et Litterman).
\item D.~Goldfarb et G.~Iyengar, ``Robust portfolio selection problems'', \emph{Mathematics of Operations Research} 28(1), 2003.
\item S.~Maillard, T.~Roncalli et J.~Teïletche, ``The properties of equally weighted risk contribution portfolios'', \emph{Journal of Portfolio Management}
36(4), 2010.
\item M.~López de Prado, ``Building diversified portfolios that outperform out of sample'', \emph{Journal of Portfolio Management} 42(4), 2016.
\item N.~Gârleanu et L.~H.~Pedersen, ``Dynamic trading with predictable returns and transaction costs'', \emph{Journal of Finance} 68(6), 2013.
\item S.~Boyd et al., ``Multi-period trading via convex optimization'', \emph{Foundations and Trends in Optimization} 3(1), 2017.
\item C.~S.~Asness, A.~Frazzini et L.~H.~Pedersen, ``Leverage aversion and risk parity'', \emph{Financial Analysts Journal} 68(1), 2012.
\end{omsources}

\section{Exercices}

\begin{exercise}[$\star$]\label{exo:rs:portfolio-construction-ii:1}
Deux actifs non corrélés ont des volatilités de 20~\% et 10~\%. Calculer la volatilité du portefeuille à pondérations égales, et les pondérations du \omterm{def:rs:portfolio-construction-ii:erc}{portefeuille à contributions au risque égales}.
\end{exercise}

\begin{exercise}[$\star$]\label{exo:rs:portfolio-construction-ii:2}
Montrer qu'en l'absence de vues, le livre de Black--Litterman est le portefeuille de marché.
\end{exercise}

\begin{exercise}[$\star$]\label{exo:rs:portfolio-construction-ii:3}
Avec $a/\gamma = 0.78$, donner les pondérations du \omterm{def:rs:portfolio-construction-ii:aim}{portefeuille cible} sur des signaux dont la décroissance vaut $\phi = 1$, $1/60$ et $0.00137$ par jour.
\end{exercise}

\begin{exercise}[$\star\star$]\label{exo:rs:portfolio-construction-ii:4}
Pourquoi l'optimisation robuste réduit-elle la plus forte réponse à une variation d'un point de base, mais non la réponse moyenne~?
\end{exercise}

\begin{exercise}[$\star\star$]\label{exo:rs:portfolio-construction-ii:5}
Pourquoi la volatilité du \omterm{def:rs:portfolio-construction-ii:erc}{portefeuille à contributions au risque égales} se situe-t-elle entre celles du portefeuille de variance minimale et du portefeuille à pondérations égales~?
\end{exercise}

\begin{exercise}[$\star\star$]\label{exo:rs:portfolio-construction-ii:6}
Pourquoi l'ajustement partiel vers le portefeuille de Markowitz donne-t-il de moins bons résultats après coûts que la règle de \omterm{def:rs:portfolio-construction-ii:aim}{portefeuille cible} exécutée au même rythme~?
\end{exercise}

\begin{exercise}[$\star\star\star$]\label{exo:rs:portfolio-construction-ii:7}
\emph{Programmation.} Exécuter \texttt{rs\_allocation.multiperiod} pour les quatre règles avec $\lambda = 1\,000$ et $10\,000$, puis expliquer quelle règle survit et pourquoi.
\end{exercise}

\begin{exercise}[$\star\star\star$]\label{exo:rs:portfolio-construction-ii:8}
\emph{Trouver la faille.} «~La \omterm{def:rs:portfolio-construction-ii:hrp}{parité de risque hiérarchique} a dépassé l'optimisation moyenne-variance de 0.10 de ratio de Sharpe après coûts sur huit ans, nous allons donc changer de méthode.~»
\end{exercise}

\section{Problème~: le retournement d'un point de base}

\begin{problem}[{Problème du week-end --- livres stables et bonnes transactions}]\label{pb:rs:portfolio-construction-ii:1}
Les livres à cinquante noms du chapitre et les règles multipériodes quotidiennes.

\medskip\noindent\textbf{Partie I --- Stabilité.}
\begin{enumerate}
\item Dans quelle mesure un point de base sur un alpha déplace-t-il chaque livre fondé sur les alphas, en moyenne et au maximum~?
\item Dans quelle mesure une prévision retirée déplace-t-elle ces livres~?
\item Pourquoi les livres fondés sur le risque ne bougent-ils pas, et pourquoi n'est-ce pas un mérite~?
\item Que retire l'ellipsoïde robuste, et que ne retire-t-il pas~?
\end{enumerate}

\medskip\noindent\textbf{Partie II --- Équilibre et risque.}
\begin{enumerate}[resume]
\item Écrire la moyenne a posteriori de Black--Litterman et expliquer $\tau$.
\item Quelles sont les volatilités et les plus fortes pondérations des livres à \omterm{def:rs:portfolio-construction-ii:erc}{contributions au risque} égales et HRP~?
\item Comment HRP évite-t-il d'inverser la matrice de covariance~?
\item Pourquoi la \omterm{def:rs:portfolio-construction-ii:erc}{parité de risque} entre classes d'actifs peut-elle rapporter plus que ne le suggère son risque~?
\end{enumerate}

\medskip\noindent\textbf{Partie III --- Performance.}
\begin{enumerate}[resume]
\item Donner les ratios de Sharpe des huit règles après coûts, ainsi que leur erreur type.
\item Que peut-on conclure de la comparaison, et que ne peut-on pas conclure~?
\end{enumerate}

\medskip\noindent\textbf{Partie IV --- Trading dans le temps.}
\begin{enumerate}[resume]
\item Quels sont les trois signaux et leurs taux de décroissance~?
\item Pourquoi la reoptimisation quotidienne perd-elle après coûts~?
\item Que pondère le \omterm{def:rs:portfolio-construction-ii:aim}{portefeuille cible}, et à quelle vitesse trade-t-il pour $\lambda = 100$~?
\item Énoncer le \emph{résultat nommé}~: le taux de rotation par point de base de perturbation de l'alpha pour chaque méthode, ainsi que le ratio de Sharpe net du \omterm{def:rs:portfolio-construction-ii:aim}{portefeuille cible} face à la reoptimisation quotidienne.
re-optimisation.
\item Comment les règles se comportent-elles lorsque les coûts sont multipliés par dix puis par cent~?
\item Pourquoi le rééquilibrage mensuel réussit-il à faibles coûts et échoue-t-il à coûts élevés~?
\item Que changeraient des prévisions bruitées~?
\item Laquelle des règles du chapitre utiliseriez-vous pour un livre de signaux quotidiens~?
\item Comment estimeriez-vous $\lambda$~?
\item En une phrase~: qu'est-ce qui compte le plus à l'horizon quotidien, le livre ou le chemin pour y parvenir~?
\end{enumerate}
\end{problem}

\section{Questions d'entretien}

\begin{interviewq}[$\star$ \iqroles{chercheur}]\label{iq:rs:portfolio-construction-ii:1}
Quel problème Black--Litterman résout-il, et comment~?
\end{interviewq}

\begin{interviewq}[$\star\star$ \iqroles{chercheur, risque}]\label{iq:rs:portfolio-construction-ii:2}
Expliquer la \omterm{def:rs:portfolio-construction-ii:erc}{parité de risque}. Que suppose-t-elle, et que lui faut-il pour offrir des rendements comparables à ceux des actions~?
\end{interviewq}

\begin{interviewq}[$\star\star$ \iqroles{chercheur}]\label{iq:rs:portfolio-construction-ii:3}
Comment l'optimisation robuste modifie-t-elle un portefeuille moyenne-variance~? Comment choisiriez-vous la taille de l'\omterm{def:rs:portfolio-construction-ii:robust}{ensemble d'incertitude}~?
\end{interviewq}

\begin{interviewq}[$\star\star$ \iqroles{chercheur}]\label{iq:rs:portfolio-construction-ii:4}
Décrire la \omterm{def:rs:portfolio-construction-ii:hrp}{parité de risque hiérarchique} et les cas où vous l'utiliseriez.
\end{interviewq}

\begin{interviewq}[$\star\star$ \iqroles{chercheur, trader}]\label{iq:rs:portfolio-construction-ii:5}
Vous disposez d'un signal rapide et d'un signal lent, ainsi que de coûts de transaction quadratiques. Comment les tradez-vous~?
\end{interviewq}

\begin{interviewq}[$\star\star\star$ \iqroles{chercheur}]\label{iq:rs:portfolio-construction-ii:6}
Calculer les pondérations du \omterm{def:rs:portfolio-construction-ii:erc}{portefeuille à contributions au risque égales} pour deux actifs, et montrer qu'avec des actifs non corrélés ce sont des pondérations inverses à la volatilité.
\end{interviewq}
